\subsection{Power Line Carrier Band}
\label{sec:plc}

\subsubsection{Sources and coupling paths}
Converter noise in the PLC band originates in the switching of the valves. In
a line-commutated converter the commutation notches excite the stray
capacitances and inductances of the valve hall and buswork into oscillations
in the range of tens to hundreds of kilohertz; in a modular multilevel
converter the individual submodule switchings are small steps, and the
emission in this band is lower by one to two orders of magnitude, but it is
not zero and still depends on the switching frequency of the control and on
the layout of the valve hall. The noise leaves the station on three paths: as
a conducted current on the AC lines through the converter transformers and
the AC switchyard, as a conducted current on the DC line, and by inductive
coupling from the DC line into any transmission line that runs parallel to or
crosses it. The third path is the one a station-terminal measurement does not
capture.

\subsubsection{Form of the requirement}
A PLC clause has two parts. The first is an absolute limit on the noise
injected into each connected line, stated as a noise voltage in a 3\,kHz
bandwidth, the channel bandwidth of a carrier receiver, measured with a
quasi-peak detector according to CISPR~16-1-1~[32] at the line side of the
switchyard. The second is a signal-to-noise criterion for each carrier
frequency in service on a protected line: the noise coupled into that line
may not reduce the carrier signal-to-noise ratio at the receiver below a
stated margin, with the minimum received carrier level and the receiver
bandwidth given by the line owner and the coupling calculated from the
geometry of the parallel or crossing section. The second part governs
wherever it is stricter. The list of protected lines and carrier frequencies
is compiled by the owner from the UTC database~[28] and from the owners of the
lines within a stated corridor of the route, and is handed to the supplier
with the DC line data. The supplier's response is the PLC filter on the AC
and DC side, tuned to the protected frequencies, and a demonstration that the
limits hold with any redundant filter branch out of service and, in a dual
symmetrical monopole, with one converter out of service.

\subsubsection{Indicative values}
A margin of 20\,dB on the carrier signal-to-noise ratio is the value in common
use for teleprotection channels. For the absolute limit, 60\,dB($\mu$V),
equal to 1\,mV in 3\,kHz, is representative of the values specified for
line-commutated schemes~[29]; whether it protects a given carrier channel
follows from the signal-to-noise criterion, which is why that criterion
governs. A VSC station meets it with a wide margin; the clause is nevertheless kept so
that the supplier must demonstrate it. A corridor of 2\,km for the survey of
protected lines covers the distance within which inductive coupling from a
DC line is calculable above the natural noise of the coupled line. For a
private network that is not connected to the wider AC grid and runs its
teleprotection over fibre, the PLC clause protects third-party lines through
the DC route only, and the AC-side part reserves carrier capability for a
later buildout.
\FloatBarrier
